%%%%%%%%%%%%%%%%%%%%%%%%%%%%%%%%%%%%%%%%%%%%%%%%%%%%%%%%%%%%%%%%%%%%%%%%
% Beamer Presentation - LaTeX - Template Version 1.0 (10/11/12)
% This template has been downloaded from: http://www.LaTeXTemplates.com
% License: % CC BY-NC-SA 3.0 (http://creativecommons.org/)
% Modified by Rahmat M. Samik-Ibrahim

% REV003 Mon 14 Feb 2022 21:50:54 WIB
% REV001 Mon 07 Feb 2022 05:54:47 WIB
% STARTX Wed 14 Sep 2016 10:45:19 WIB
%%%%%%%%%%%%%%%%%%%%%%%%%%%%%%%%%%%%%%%%%%%%%%%%%%%%%%%%%%%%%%%%%%%%%%%%%

% PACKAGES AND THEMES 
\documentclass[aspectratio=169, xcolor=table, notheorems, hyperref={pdfpagelabels=false}]{beamer}
\input{beamer.tex}
\newcommand{\revision}{REV021 06-Feb-2023}
% w! tmptmp
% REV021: Mon 06 Feb 2023 00:00
% REV019: Thu 02 Feb 2023 19:00
% REV009: Mon 18 Apr 2022 06:00
% REV007: Thu 17 Mar 2022 10:00
% REV001: Mon 07 Feb 2022 18:00
% STARTS: Wed 24 Aug 2016 19:00
%%%%%%%%%%%%%%%%%%%%%%%%%%%%%%%
\newcommand{\kopikopi}{\textcopyright{}2016-2023 CBKadal + VauLSMorg}



% XXXXXXXXXXXXXXXXXXXXXXXXXXXXXXXXXXXXXXXXXXXXXXXXXXXXXXXXXXXXXXXXXXXXXXXXXX
% The short title appears at the bottom of every slide, 
% the full title is only on the title page
% \date{\today}
\title[\kopikopi]
{CSCE604227 System Programming \\
CSCE604227 Pemrograman Sistem \\
Week 01:
Linux Kernel and Programming Interface}
\author{C. BinKadal}
\institute[SDN]
{
Sendirian Berhad\\
\medskip
\url{https://sp.vlsm.org/Slides/sp01.pdf}
\\ \texttt{Always check for the latest revision!}
}
\date{\revision}

% XXXXXXXXXXXXXXXXXXXXXXXXXXXXXXXXXXXXXXXXXXXXXXXXXXXXXXXXXXXXXXXXXXXXXXXXXX
\begin{document}

\lstset{
basicstyle=\ttfamily\tiny, % \tiny \small \footnotesize 
breakatwhitespace=true,
language=C,
columns=fullflexible,
keepspaces=true,
breaklines=true,
tabsize=3, 
showstringspaces=false,
extendedchars=true}

\section{Start}
\begin{frame}
\titlepage
\end{frame}

% XXXXXXXXXXXXXXXXXXXXXXXXXXXXXXXXXXXXXXXXXXXXXXXXXXXXXXXXXXXXXXXXXXXXXXXXXX

%%%%%%%%%%%%%%%%%%%%%%%%%%%%%%%
% REV019: Thu 02 Feb 2023 00:00
% REV003: Mon 14 Feb 2022 21:00
% REV001: Mon 07 Feb 2022 18:00
% START0: Wed 14 Sep 2016 10:00
%%%%%%%%%%%%%%%%%%%%%%%%%%%%%%%

\begin{frame}[fragile]
\section{Schedule}
\frametitle{SP221\footnote{%
) This information will be on \textbf{EVERY} page two (2) of this course material.}): 
System Progamming}

\vspace{5pt}

\scalebox{1.2}{%
\begin{tabular}{|c|c|l|}
\hline
\textbf{Week} & 
\textbf{Topic} \\ 
\hline
Week 00  & Overview \\
Week 01  & Linux Kernel and Programming Interface \\
Week 02  & Revisit Linux From Scratch \\
Week 03  & FUSE: Filesystem in Userspace  \\
Week 04  & Project FUSE 1 \\
Week 05  & Project FUSE 2 \\
\hline
Week 06  & Project FUSE 3 \\
Week 07  & Project FUSE 4 \\
Week 08  & Project FUSE 5 \\
Week 09  & Project FUSE 6 \\
Week 10  & Project FUSE Presentation \\
% MidTerm  & 00 XXX 2020 (XX:XX-XX:XX) & MidTerm (UTS) & \cellcolor{red!44} TBA! \\
% Reserved & 00 XXX - 00 XXX 2020 & Q \& A & \\
% Final    & 00 XXX 2020 XX:XX & First Part Final  (UAS tahap I)  & \cellcolor{red!44} This schedule is   \\
% Extra    & NA & No Extra assignment & \cellcolor{red!44} subject to change. \\
\hline
\end{tabular}
}
\end{frame}

\begin{frame}[fragile]
\frametitle{\textbf{STARTING POINT} --- 
{
\definecolor{links}{HTML}{FDEE00}
\hypersetup{colorlinks,linkcolor=,urlcolor=links}
\url{https://sp.vlsm.org/}
}
}
\begin{itemize}
\item[$\square$] \textbf{Text Book} ---
The Linux Programming Interface, 2010, No Starch Press, 
ISBN 978-1-59327-220-3 --- \url{https://man7.org/tlpi/}. 
\item[$\square$] \textbf{Resources}
\begin{itemize}
\item[$\square$] \href{https://scele.cs.ui.ac.id/course/view.php?id=3545}{\textbf{SCELE}} ---
\url{https://scele.cs.ui.ac.id/course/view.php?id=3545}.\\
The enrollment key is \textbf{XXX}.
\item[$\square$] \textbf{Download Slides and Demos from GitHub.com} \\
\url{https://github.com/os2xx/sysprog/}:

                 {\scriptsize%
                 \href{https://sp.vlsm.org/Slides/sp00.pdf}{\texttt{sp00.pdf} (W00)},
                 \href{https://sp.vlsm.org/Slides/sp01.pdf}{\texttt{sp01.pdf} (W01)},
                 \href{https://sp.vlsm.org/Slides/sp02.pdf}{\texttt{sp02.pdf} (W02)},
                 \href{https://sp.vlsm.org/Slides/sp03.pdf}{\texttt{sp03.pdf} (W03)},

                 \href{https://sp.vlsm.org/Slides/sp04.pdf}{\texttt{sp04.pdf} (W04)},
                 \href{https://sp.vlsm.org/Slides/sp05.pdf}{\texttt{sp05.pdf} (W05)},
                 \href{https://sp.vlsm.org/Slides/sp06.pdf}{\texttt{sp06.pdf} (W06)},
                 \href{https://sp.vlsm.org/Slides/sp07.pdf}{\texttt{sp07.pdf} (W07)},

                 \href{https://sp.vlsm.org/Slides/sp08.pdf}{\texttt{sp08.pdf} (W08)},
                 \href{https://sp.vlsm.org/Slides/sp09.pdf}{\texttt{sp09.pdf} (W09)},
                 \href{https://sp.vlsm.org/Slides/sp10.pdf}{\texttt{sp10.pdf} (W10)}.
                 }
\item[$\square$] \textbf{LFS} --- \url{http://www.linuxfromscratch.org/lfs/view/stable/}
\item[$\square$] \textbf{OSP4DISS} --- \url{https://osp4diss.vlsm.org/}
\item[$\square$] \textbf{This is How Me DO IT!} --- \url{https://doit.vlsm.org/}
\begin{itemize}
\item[$\square$] PS: ''Me'' rhymes better than ''I'' duh!
\end{itemize}
\end{itemize}
\end{itemize}
\end{frame}



% XXXXXXXXXXXXXXXXXXXXXXXXXXXXXXXXXXXXXXXXXXXXXXXXXXXXXXXXXXXXXXXXXXXXXXXXXX
% Throughout your presentation, if you choose to use \section{} and 
% \subsection{} commands, these will automatically be printed on 
% this slide as an overview of your presentation
\section{Agenda}
\begin{frame}{Outline}
  \frametitle{Agenda}
  \tableofcontents[sections={1-}]
\end{frame}
% \begin{frame}
%    \frametitle{Agenda (2)}
%    \tableofcontents[sections={12-}]
% \end{frame}

% XXXXXXXXXXXXXXXXXXXXXXXXXXXXXXXXXXXXXXXXXXXXXXXXXXXXXXXXXXXXXXXXXXXXXXX
\section{The Linux Kernel Archives}
\begin{frame}[fragile]
\frametitle{The Linux Kernel Archives}
\begin{itemize}
\item URL: \url{https://kernel.org/}
\begin{itemize}
\item HTTP: \url{https://www.kernel.org/pub/}
\item Kernel Source: \url{https://www.kernel.org/pub/linux/kernel/}
\end{itemize}
\item 14-Feb-2022 --- 5.16.9 --- \url{https://www.kernel.org/finger_banner}
% \begin{lstlisting}[basicstyle=\ttfamily\tiny]
\begin{lstlisting}[basicstyle=\ttfamily\small]
The latest stable version of the Linux kernel is:             5.16.9
The latest mainline version of the Linux kernel is:           5.17-rc4
The latest stable 5.16 version of the Linux kernel is:        5.16.9
The latest longterm 5.15 version of the Linux kernel is:      5.15.23
The latest longterm 5.10 version of the Linux kernel is:      5.10.100
[...]
The latest longterm 4.9 version of the Linux kernel is:       4.9.301
The latest longterm 4.4 version of the Linux kernel is:       4.4.302 (EOL)
The latest linux-next version of the Linux kernel is:         next-20220214
\end{lstlisting}
\item How to compile Linux Kernel on Debian VirtualBox Guest --- \url{https://doit.vlsm.org/007.html}.
\item Debian Packages List: \url{https://osp4diss.vlsm.org/osp-103.html}
\end{itemize}
\end{frame}

% XXXXXXXXXXXXXXXXXXXXXXXXXXXXXXXXXXXXXXXXXXXXXXXXXXXXXXXXXXXXXXXXXXXXXXX
\section{The Linux Programming Interface}
\begin{frame}[fragile]
\frametitle{The Linux Programming Interface (1)}
\begin{itemize}
\item URL: \url{https://man7.org/tlpi/}
\item Purchasing: \url{https://man7.org/tlpi/purchase.html}
\item Source Code: \url{https://man7.org/tlpi/code/}
\begin{itemize}
\item Lastest Distribution: \url{https://man7.org/tlpi/code/download/tlpi-220211-dist.tar.gz}
\item Download Full List: \url{https://man7.org/tlpi/code/download/}
\end{itemize}
\item README: \url{https://man7.org/tlpi/code/README.html}
\item BUILDING notes: \url{https://man7.org/tlpi/code/BUILDING.html}
\item FAQ: \url{https://man7.org/tlpi/code/faq.html}
\item Debian Packages List: \url{https://osp4diss.vlsm.org/osp-103.html}
\end{itemize}
\end{frame}

% XXXXXXXXXXXXXXXXXXXXXXXXXXXXXXXXXXXXXXXXXXXXXXXXXXXXXXXXXXXXXXXXXXXXXXX
\begin{frame}[fragile]
\frametitle{The Linux Programming Interface (2)}
\begin{itemize}
\item Test: \texttt{time/calendar\_time}

% \begin{lstlisting}[basicstyle=\ttfamily\tiny]
\begin{lstlisting}[basicstyle=\ttfamily\small]

cbkadal@cbkadal:~/src/tlpi-dist$ time/calendar_time 
Seconds since the Epoch (1 Jan 1970): 1644855271 (about 52.123 years)
  gettimeofday() returned 1644855271 secs, 35730 microsecs
Broken down by gmtime():
  year=122 mon=1 mday=14 hour=16 min=14 sec=31 wday=1 yday=44 isdst=0
Broken down by localtime():
  year=122 mon=1 mday=14 hour=23 min=14 sec=31 wday=1 yday=44 isdst=0

asctime() formats the gmtime() value as: Mon Feb 14 16:14:31 2022
ctime() formats the time() value as:     Mon Feb 14 23:14:31 2022
mktime() of gmtime() value:    1644830071 secs
mktime() of localtime() value: 1644855271 secs

cbkadal@cbkadal:~/src/tlpi-dist$

\end{lstlisting}
\end{itemize}
\end{frame}

% XXXXXXXXXXXXXXXXXXXXXXXXXXXXXXXXXXXXXXXXXXXXXXXXXXXXXXXXXXXXXXXXXXXXXXX
\begin{frame}[fragile]
\frametitle{The Linux Programming Interface (3)}
\begin{itemize}
\item Test: \texttt{getopt/t\_getopt} (TLPI Appendix B)
% \begin{lstlisting}[basicstyle=\ttfamily\small]
\begin{lstlisting}[basicstyle=\ttfamily\tiny]
cbkadal@cbkadal:~/src/tlpi-dist$ getopt/t_getopt -x -p hello world
opt =120 (x); optind = 2
opt =112 (p); optind = 4
-x was specified (count=1)
-p was specified with the value "hello"
First nonoption argument is "world" at argv[4]

cbkadal@cbkadal:~/src/tlpi-dist$ getopt/t_getopt -p
opt = 58 (:); optind = 2; optopt =112 (p)
Missing argument (-p)
Usage: getopt/t_getopt [-p arg] [-x]

cbkadal@cbkadal:~/src/tlpi-dist$ getopt/t_getopt -a
opt = 63 (?); optind = 2; optopt = 97 (a)
Unrecognized option (-a)
Usage: getopt/t_getopt [-p arg] [-x]

cbkadal@cbkadal:~/src/tlpi-dist$ getopt/t_getopt -p str -- -x
opt =112 (p); optind = 3
-p was specified with the value "str"
First nonoption argument is "-x" at argv[4]

cbkadal@cbkadal:~/src/tlpi-dist$ getopt/t_getopt -p -x
opt =112 (p); optind = 3
-p was specified with the value "-x"

\end{lstlisting}
\end{itemize}
\end{frame}

\end{document}

